\documentclass[12pt,a4paper]{article}

\usepackage[left=17mm,right=10mm,top=20mm,bottom=20mm]{geometry}
\usepackage{iftex}
\ifPDFTeX
  \usepackage[T2A]{fontenc}
  \usepackage[utf8]{inputenc}
  \usepackage[russian,english]{babel}
\else
  \usepackage{fontspec}
  \usepackage{polyglossia}
  \setdefaultlanguage{russian}
  \setotherlanguage{english}
  \IfFontExistsTF{Times New Roman}{
    \setmainfont{Times New Roman}
  }{
    \setmainfont{TeX Gyre Termes}
  }
  \IfFontExistsTF{Times New Roman}{
    \setmonofont{Times New Roman}
  }{
    \setmonofont{TeX Gyre Termes}
  }
\fi
\usepackage{csquotes}

\usepackage{setspace}
\onehalfspacing
\usepackage{indentfirst}
\setlength{\parindent}{1.25cm}
\setlength{\parskip}{0pt}
\usepackage{microtype}
\usepackage{hyperref}
\usepackage{cite}
\usepackage{fancyhdr}
\usepackage{float}
\usepackage{array}
\usepackage{graphicx}
\usepackage{caption}
\usepackage{longtable}
\usepackage{booktabs}
\usepackage{xcolor}
\usepackage{tikz}
\usetikzlibrary{positioning,arrows.meta,shapes.geometric}

\captionsetup[table]{labelfont=bf,textfont=it}
\captionsetup[figure]{labelfont=bf,textfont=it}

\hypersetup{
  colorlinks=true,
  linkcolor=blue,
  citecolor=blue,
  urlcolor=blue
}

\renewcommand{\contentsname}{Содержание}
\renewcommand{\refname}{Список использованных источников}
\newcolumntype{L}[1]{>{\raggedright\arraybackslash}p{#1}}
\newcommand{\code}[1]{{\small\texttt{#1}}}

\makeatletter
\newcommand{\papercenterpagenumber}{%
  \if@twoside
    \ifodd\c@page
      \hspace*{-\dimexpr 1in + \hoffset + \oddsidemargin\relax}%
    \else
      \hspace*{-\dimexpr 1in + \hoffset + \evensidemargin\relax}%
    \fi
  \else
    \hspace*{-\dimexpr 1in + \hoffset + \oddsidemargin\relax}%
  \fi
  \makebox[\paperwidth]{\hfil\thepage\hfil}%
}
\makeatother

\fancypagestyle{plain}{%
  \fancyhf{}%
  \renewcommand{\headrulewidth}{0pt}%
  \renewcommand{\footrulewidth}{0pt}%
  \fancyfoot[L]{\makebox[0pt][l]{\papercenterpagenumber}}%
}

\begin{document}

\IfFileExists{title_kr.tex}{
  % Титульный лист для итогового отчёта (шаблон ФКН)
\begin{titlepage}
\thispagestyle{empty}

{\setstretch{1.0}
\begin{center}
ПРАВИТЕЛЬСТВО РОССИЙСКОЙ ФЕДЕРАЦИИ\\
ФГАОУ ВО НАЦИОНАЛЬНЫЙ ИССЛЕДОВАТЕЛЬСКИЙ УНИВЕРСИТЕТ\\
«ВЫСШАЯ ШКОЛА ЭКОНОМИКИ»
\\
\bigskip
Факультет компьютерных наук\\
Образовательная программа «Прикладная математика и информатика»
\end{center}
}

\vspace{2em}
% УДК нужно указывать только для исследовательского проекта
% УДК ХХХХХ
\vspace{4em}

\begin{center}
{\bf Отчет о программном проекте на тему:}\\
{\bf DogPaw: мобильная платформа для владельцев собак}\\
\end{center}

\vspace{3em}

\noindent{\bf Выполнил студент:}

\vspace{2mm}
{\setstretch{1.15}
\noindent\begin{tabular}{@{}p{0.42\textwidth}@{\hspace{1.5em}}p{0.48\textwidth}@{}}
группы \#БПМИ237, 3 курса & Диалектов Иван Фёдорович \\
\end{tabular}}

\vspace{1.5em}
\noindent{\bf Приняли руководители проекта:}

\vspace{2mm}
{\setstretch{1.15}
\begin{tabular}{@{}p{0.92\textwidth}@{}}
Песоцкая Елена Юрьевна \\
Кандидат экономических наук \\
Доцент \\
Национальный исследовательский университет «Высшая школа экономики» (НИУ ВШЭ), \\
Факультет компьютерных наук, Департамент программной инженерии \\
\end{tabular}}

\vspace{\fill}

\begin{center}
Москва 2026
\end{center}

\end{titlepage}

}{
  \begin{titlepage}
  \thispagestyle{empty}
  {\setstretch{1.0}
  \begin{center}
  ПРАВИТЕЛЬСТВО РОССИЙСКОЙ ФЕДЕРАЦИИ\\
  ФГАОУ ВО НАЦИОНАЛЬНЫЙ ИССЛЕДОВАТЕЛЬСКИЙ УНИВЕРСИТЕТ\\
  «ВЫСШАЯ ШКОЛА ЭКОНОМИКИ»
  \\
  \bigskip
  Факультет компьютерных наук\\
  Образовательная программа «Прикладная математика и информатика»
  \end{center}
  }
  \vspace{2em}
  \vspace{4em}
  \begin{center}
  {\bf Отчет о программном проекте на тему:}\\
  {\bf DogPaw: мобильная платформа для владельцев собак}\\
  \end{center}
  \vspace{3em}
  \noindent{\bf Выполнил студент:}
  \vspace{2mm}
  {\setstretch{1.15}
  \noindent\begin{tabular}{@{}p{0.42\textwidth}@{\hspace{1.5em}}p{0.48\textwidth}@{}}
  группы \#БПМИ237, 3 курса & Диалектов Иван Фёдорович \\
  \end{tabular}}
  \vspace{1.5em}
  \noindent{\bf Приняли руководители проекта:}
  \vspace{2mm}
  {\setstretch{1.15}
  \begin{tabular}{@{}p{0.92\textwidth}@{}}
  Песоцкая Елена Юрьевна \\
  Кандидат экономических наук \\
  Доцент \\
  Национальный исследовательский университет «Высшая школа экономики» (НИУ ВШЭ), \\
  Факультет компьютерных наук, Департамент программной инженерии \\
  \end{tabular}}
  \vspace{\fill}
  \begin{center}
  Москва 2026
  \end{center}
  \end{titlepage}
}

\pagenumbering{arabic}
\setcounter{page}{2}
\pagestyle{plain}

\tableofcontents
\clearpage

\section*{Аннотация}
\addcontentsline{toc}{section}{Аннотация}

В отчете представлены результаты финального этапа программного проекта \textbf{DogPaw} --- мобильной платформы для владельцев собак с клиент-серверной архитектурой. Цель проекта заключалась в создании единого цифрового продукта, который объединяет учет данных о питомце, трекинг прогулок, социальные функции, карту dog-friendly мест и инструменты безопасности (QR-идентификация и режим поиска потерявшегося питомца). 

В результате реализована рабочая версия приложения для \textit{iOS/Android} на \textit{React Native (Expo)} и серверной части на \textit{Go + Gin + GORM + SQLite}. На финальном этапе были добавлены: регистрация и вход по JWT, персональные и ролевые права, лента с фото/лайками/комментариями и сортировкой, добавление друзей по email и по QR, модерация пользовательских статей для базы знаний, карта с фильтрами категорий и логикой приватности геопозиции. Дополнительно доведены до актуального состояния: трекер прогулок (время с секундами, история и удаление записей), поиск мест на backend по данным OpenStreetMap (Overpass) с кэшем, режим \textit{lost\&find} с текстом объявления и снятием метки с карты, отдельный экран просмотра статьи в базе знаний.

Проведено тестирование функциональности на уровне API и клиентской логики, выполнены повторные прогоны после исправлений, оформлен протокол тестирования. Показано, что система поддерживает ключевые сценарии предметной области, а архитектурные решения обеспечивают расширяемость проекта для дальнейшего перехода к production-сценарию.

\section*{Ключевые слова}
\addcontentsline{toc}{section}{Ключевые слова}
мобильное приложение; pet-tech; React Native; Expo; Golang; REST API; JWT; геолокация; социальная лента; модерация контента.


\section{Введение}

\subsection{Актуальность темы}

Сегмент pet-tech демонстрирует устойчивый рост: владельцы животных все чаще используют мобильные сервисы для учета здоровья питомца, поиска профильных услуг и коммуникации с сообществом~\cite{statista_pet}. При этом повседневный пользовательский путь обычно фрагментирован между несколькими приложениями: одно для заметок о питомце, другое для карт, третье для общения. Такая фрагментация снижает качество пользовательского опыта, увеличивает число ручных действий и мешает построению единой истории данных.

Для владельцев собак эта задача особенно важна по двум причинам. Во-первых, приложение используется в регулярных бытовых сценариях: прогулки, контроль состояния питомца и общение с другими владельцами. Во-вторых, есть критичные по времени ситуации (потеря питомца, быстрый поиск контактов владельца), где требуется надежная и быстро доступная функциональность.

\subsection{Проблема и цель проекта}

Проблема проекта формулируется как отсутствие доступного единого мобильного решения, которое одновременно покрывает: профиль питомца, прогулочную активность, карту полезных мест, социальный блок и контентную базу знаний с проверкой качества материалов.

\textbf{Цель проекта} --- разработать и довести до работоспособного состояния программную платформу DogPaw, обеспечивающую комплексный набор пользовательских сценариев владельца собаки в одной системе.

\subsection{Задачи финального этапа}

На финальном этапе были поставлены и решены следующие задачи:
\begin{enumerate}
  \item завершить серверную модель авторизации и прав доступа;
  \item реализовать полнофункциональную ленту: публикации с фото, лайки, комментарии, сортировка;
  \item реализовать рабочие сценарии дружбы: добавление по email/QR и удаление;
  \item доработать умную карту: фильтрация категорий, правила видимости, обновление геопозиции;
  \item ввести ролевую модель администрирования и модерацию базы знаний;
  \item привести интерфейс ключевых экранов к единой визуальной концепции;
  \item провести системное тестирование и оформить результаты в технической документации.
\end{enumerate}

\subsection{Объект и предмет разработки}

Объект исследования --- мобильные информационные системы для владельцев домашних животных.  
Предмет разработки --- программная архитектура и реализация клиент-серверного приложения DogPaw, включая логику данных, пользовательские сценарии и правила модерации контента.


\section{Анализ предметной области и аналогов}

\subsection{Пользовательские потребности}

На основе анализа пользовательских сценариев и существующих решений можно выделить пять базовых потребностей:
\begin{itemize}
  \item хранение и быстрый доступ к актуальным данным о питомце;
  \item безопасная идентификация питомца и контакт с владельцем;
  \item понятный учет прогулочной активности;
  \item поиск релевантных мест и взаимодействие с локальным сообществом;
  \item доступ к полезному контенту с контролем качества.
\end{itemize}

\subsection{Сравнение с аналогами}

Для корректной постановки задачи сравнение выполнялось не с абстрактными \enquote{типовыми} приложениями, а с реальными прямыми и косвенными аналогами:
\begin{itemize}
  \item \textbf{PetDesk} --- профиль питомца, напоминания, коммуникация с клиникой~\cite{petdesk};
  \item \textbf{Rover} --- услуги догситтинга и выгулов, marketplace-модель~\cite{rover};
  \item \textbf{Strava} --- зрелая модель трекинга активности и социальной ленты~\cite{strava};
  \item \textbf{BringFido} --- каталог dog-friendly мест~\cite{bringfido};
  \item \textbf{PawBoost} --- сценарии lost\&found для домашних животных~\cite{pawboost}.
\end{itemize}

Тезис \enquote{отсутствует единое решение} в проекте формализован так: среди рассмотренных продуктов не обнаружено приложения, которое одновременно предоставляет \textit{в одном пользовательском контуре} (без переключения между сервисами) следующие функции: профиль питомца + QR-идентификация + трекинг прогулок + карта мест по собачьим категориям + realtime-карта друзей + встроенная лента + база знаний с модерацией.

\subsection{Функциональная матрица}

Ниже приведена расширенная сравнительная матрица функций по ключевым конкурентам (та же структура, что в отчёте КТ-1). Источник: анализ конкурентов от 23.12.2025. Легенда: \textbf{+}~--- функция есть; \textbf{---}~--- нет (в том числе при частичной или иной форме реализации; для единообразия без шкалы «частично»). Сокращения колонок: Проф.\ = профиль питомца; QR = QR-адресник/скан; Прог.\ = трекер прогулок; Корм.\ = кормление (калькулятор/напоминания); Карта = карта мест; Отз.\ = оценки/отзывы; Друз.\ = друзья/комьюнити; Чаты = чаты; Лента = лента фото/видео; Live = realtime-карта людей/статусы; База = база знаний; L\&F = Lost\&Found; Серв.\ = сервисы/маркетплейс.

\footnotesize
\setlength{\tabcolsep}{2pt}
\setlength\LTleft{\fill}
\setlength\LTright{\fill}
\begin{longtable}{@{}p{1.55cm}p{1.45cm}ccccccccccccc@{}}
\caption{Матрица функций: DogPaw и конкуренты (приложение для владельцев собак)} \label{tab:analogs-kt2} \\
\toprule
\textbf{Приложение} & \textbf{Катег.} & Проф. & QR & Прог. & Корм. & Карта & Отз. & Друз. & Чаты & Лента & Live & База & L\&F & Серв. \\
\midrule
\endfirsthead

\multicolumn{15}{c}{\tablename\ \thetable\ --- продолжение} \\
\toprule
\textbf{Приложение} & \textbf{Катег.} & Проф. & QR & Прог. & Корм. & Карта & Отз. & Друз. & Чаты & Лента & Live & База & L\&F & Серв. \\
\midrule
\endhead

\bottomrule
\endfoot

DogPaw (реализация КТ-2) & all-in-one & + & + & + & + & -- & + & -- & + & + & + & + & -- \\
DogPack & карта, комьюнити, L\&F & -- & -- & -- & -- & + & -- & + & + & + & -- & -- & + & + \\
Pawscout & безопасность, комьюнити & + & -- & + & -- & + & -- & + & -- & -- & -- & -- & + & -- \\
BringFido & dog-friendly каталог & -- & -- & -- & -- & + & + & -- & -- & -- & -- & -- & -- & -- \\
Sniffspot & бронирование площадок & -- & -- & -- & -- & + & + & -- & -- & -- & -- & -- & -- & + \\
PawBoost & lost\&found & -- & -- & -- & -- & -- & -- & -- & -- & -- & -- & -- & + & -- \\
PetHub & QR ID, L\&F сервис & + & + & -- & -- & -- & -- & -- & -- & -- & -- & -- & + & -- \\
11pets & здоровье, дневник & + & -- & -- & -- & -- & -- & -- & -- & -- & -- & -- & -- & -- \\
DogLog & уход, ремайндеры & + & -- & -- & -- & -- & -- & -- & -- & -- & -- & -- & -- & -- \\
PetDesk & ремайндеры, клиники & + & -- & -- & -- & -- & -- & -- & -- & -- & -- & -- & -- & -- \\
Wag! & маркетплейс выгул & -- & -- & -- & -- & -- & + & -- & + & -- & -- & -- & -- & + \\
Rover & маркетплейс ситтеры & -- & -- & -- & -- & -- & + & -- & + & -- & -- & -- & -- & + \\
Fi (smart collar) & GPS ошейник, health & -- & -- & -- & -- & -- & -- & -- & -- & -- & -- & -- & -- & -- \\
Tractive & GPS трекер, здоровье & -- & -- & -- & -- & -- & -- & -- & -- & -- & -- & -- & -- & -- \\
Whistle & GPS трекер (закрыт) & -- & -- & -- & -- & -- & -- & -- & -- & -- & -- & -- & -- & -- \\
WoofTrax & трекер + благотворит. & -- & -- & + & -- & -- & -- & -- & -- & -- & -- & -- & -- & -- \\
Amiko & трекер прогулок & -- & -- & + & -- & -- & -- & -- & -- & -- & -- & -- & -- & -- \\
Pet First Aid & энциклопедия, помощь & + & -- & -- & -- & -- & -- & -- & -- & -- & -- & + & -- & -- \\
Puppr & обучение & -- & -- & -- & -- & -- & -- & -- & -- & -- & -- & + & -- & -- \\
Pupford & обучение, контент & + & -- & -- & -- & -- & -- & -- & -- & -- & -- & + & -- & -- \\
Dogo & обучение, трекинг & -- & -- & -- & -- & -- & -- & -- & -- & -- & + & -- & -- & -- \\
Petzbe & соцсеть питомцев & + & -- & -- & -- & -- & -- & -- & -- & + & -- & -- & -- & -- \\
\bottomrule
\end{longtable}
\normalsize

Из матрицы следует: у DogPaw в текущей версии нет отдельных чатов и экрана отзывов по точкам на карте (комментарии реализованы в ленте; отзывы по местам --- только потенциал API). При этом сочетание профиля, QR, прогулок, карты мест, друзей и ленты, умной карты (Live), базы знаний и lost\&find в одном приложении у аналогов встречается редко. Ближе всего по охвату --- DogPack и Pawscout, но у них не доведены до уровня столбца QR, кормление и Live-карта.


\section{Требования к системе}

\subsection{Функциональные требования}

Ключевые функциональные требования:
\begin{enumerate}
  \item регистрация/аутентификация пользователя и восстановление сессии;
  \item создание, изменение и просмотр карточки питомца;
  \item генерация QR-кода питомца и доступ к карточке по QR;
  \item запуск/остановка прогулки с сохранением статистики, просмотр истории и удаление записей;
  \item работа карты мест с фильтрацией по категориям и отображением собственной геопозиции;
  \item умная карта с геопозицией пользователя и друзей, цветовая индикация статуса на маркерах;
  \item социальный модуль: друзья, посты, комментарии, лайки к постам;
  \item база знаний: публикация предложений, просмотр статьи на отдельном экране, модерация;
  \item административные операции удаления контента и выдачи прав.
\end{enumerate}

\subsection{Нефункциональные требования}

Нефункциональные требования заданы через измеримые метрики и критерии приемки (таблица~\ref{tab:nfr}).

\begin{table}[H]
\centering
\small
\caption{Формализованные нефункциональные требования}
\label{tab:nfr}
\begin{tabular}{|L{0.19\textwidth}|L{0.30\textwidth}|L{0.24\textwidth}|L{0.20\textwidth}|}
\hline
\textbf{Категория} & \textbf{Метрика} & \textbf{Целевое значение} & \textbf{Метод измерения} \\
\hline
Производительность API & p95 latency для \code{GET /feed}, \code{/map/users}, \code{/articles} & \(\leq 700\) мс при 50 RPS & k6 / JMeter, 5 мин прогон \\
\hline
Отзывчивость UI & Time-to-interactive ключевых экранов & \(\leq 2.0\) c на среднем устройстве & React Native profiler + ручной замер \\
\hline
Надежность сети & Доля успешных запросов при нестабильной сети & \(\geq 99\%\) при retry=1; базовый timeout API 10\,с (для тяжёлого поиска мест на клиенте задан увеличенный лимит) & chaos-network тест в dev среде \\
\hline
Доступность backend & Uptime локального стенда во время сессии теста & \(\geq 99\%\) за 24 часа & health-check polling \\
\hline
Безопасность доступа & Доля защищенных endpoint, требующих JWT & 100\% приватных endpoint & ревизия роутов + интеграционные тесты \\
\hline
Консистентность данных & Ошибки сериализации/десериализации JSON моделей & 0 критических ошибок на smoke suite & API smoke + контрактные тесты \\
\hline
\end{tabular}
\end{table}

\textbf{Критерий \enquote{без визуальной деградации}} в проекте трактуется как отсутствие заметных лагов интерфейса при 60 FPS-ориентированном рендеринге: длительность блокирующих операций на основном потоке не должна превышать 100 мс чаще чем в 5\% интеракций экрана.

\subsection{Ограничения проекта}

\begin{itemize}
  \item локальная разработка и тестирование выполняются в LAN-сети;
  \item полнота выдачи точек на карте мест зависит от покрытия региона в OpenStreetMap и доступности публичных Overpass-инстансов;
  \item медиа в ленте на текущем этапе сохраняются как URI/данные запроса без отдельного CDN;
  \item отсутствует production-процедура миграций и централизованный мониторинг.
\end{itemize}


\section{Архитектура и проектные решения}

\subsection{Общая архитектура}

Система реализована как клиент-серверное приложение:
\begin{itemize}
  \item мобильный клиент: React Native + Expo;
  \item backend API: Go + Gin;
  \item хранение данных: SQLite + GORM;
  \item транспорт: HTTP/JSON;
  \item авторизация: JWT Bearer Token.
\end{itemize}

\subsection{Логическая декомпозиция}

\begin{table}[H]
\centering
\small
\caption{Декомпозиция модулей системы}
\begin{tabular}{|L{0.3\textwidth}|L{0.62\textwidth}|}
\hline
\textbf{Модуль} & \textbf{Ответственность} \\
\hline
Auth & Регистрация, вход, проверка токена, получение текущего пользователя, выдача admin-прав owner-аккаунтом \\
\hline
Pets & Профили питомцев, QR-идентификация, обновление медицинских данных \\
\hline
Walks & Учет прогулок, хранение статистики и маршрутов \\
\hline
Map/Places & Поиск и выдача мест, фильтры категорий, геопозиция пользователей \\
\hline
Friends/Social & Дружба, лента, лайки, комментарии, удаление постов админом \\
\hline
Knowledge Base & Публичные статьи, подача на модерацию, админ-решения approve/reject/delete \\
\hline
\end{tabular}
\end{table}

\subsection{Модель данных}

Ключевые сущности: \code{User}, \code{Pet}, \code{Walk}, \code{Place}, \code{UserLocation}, \code{Friendship}, \code{FeedPost}, \code{FeedLike}, \code{FeedComment}, \code{Article}, \code{LostPetAlert}.  
На финальном этапе добавлены поля для ролевой модели и модерации (\code{is\_admin}, \code{status}, \code{submitted\_by}, \code{approved\_by}).

\subsection{Ролевая модель}

\begin{itemize}
  \item \textbf{Пользователь:} базовые операции со своим контентом и взаимодействие с друзьями.
  \item \textbf{Администратор:} модерация базы знаний и удаление постов.
  \item \textbf{Owner-аккаунт:} единственный субъект, который может выдавать admin-роль другим пользователям (при условии дружбы).
\end{itemize}

Такой подход минимизирует риск бесконтрольной эскалации привилегий и соответствует требованию централизованного контроля прав.

\subsection{C4: Context}

\begin{figure}[H]
\centering
\begin{tikzpicture}[node distance=10mm, >=Stealth]
\node[draw, rounded corners, fill=blue!10, minimum width=42mm, minimum height=10mm] (dogpaw) {Система DogPaw};
\node[draw, rounded corners, fill=green!10, above left=of dogpaw, minimum width=38mm] (owner) {Владелец собаки};
\node[draw, rounded corners, fill=green!10, below left=of dogpaw, minimum width=38mm] (admin) {Администратор};
\node[draw, rounded corners, fill=orange!10, above right=of dogpaw, minimum width=45mm] (maps) {Картографический API};
\node[draw, rounded corners, fill=orange!10, below right=of dogpaw, minimum width=45mm] (expo) {Expo mobile runtime};
\draw[->] (owner) -- node[above, sloped] {\tiny мобильный UI} (dogpaw);
\draw[->] (admin) -- node[below, sloped] {\tiny модерация} (dogpaw);
\draw[->] (dogpaw) -- node[above, sloped] {\tiny поиск мест} (maps);
\draw[->] (dogpaw) -- node[below, sloped] {\tiny клиентская сборка} (expo);
\end{tikzpicture}
\caption{C4 Context: внешнее окружение системы DogPaw}
\end{figure}

\subsection{C4: Container}

\begin{figure}[H]
\centering
\begin{tikzpicture}[node distance=8mm, >=Stealth]
\node[draw, rounded corners, fill=purple!10, minimum width=38mm] (mobile) {Mobile App\\(React Native)};
\node[draw, rounded corners, fill=cyan!10, right=20mm of mobile, minimum width=38mm] (api) {Backend API\\(Go/Gin)};
\node[draw, rounded corners, fill=yellow!15, right=20mm of api, minimum width=34mm] (db) {SQLite DB};
\node[draw, rounded corners, fill=gray!15, below=12mm of api, minimum width=42mm] (ext) {External Maps Provider};
\draw[->] (mobile) -- node[above]{\tiny HTTPS/JSON} (api);
\draw[->] (api) -- node[above]{\tiny GORM} (db);
\draw[->] (api) -- node[right]{\tiny REST} (ext);
\end{tikzpicture}
\caption{C4 Container: контейнеры и их взаимодействие}
\end{figure}

\subsection{Sequence diagram (основной сценарий ленты)}

\begin{figure}[H]
\centering
\begin{tikzpicture}[node distance=9mm, >=Stealth]
\node (u) at (0,0) {\textbf{User}};
\node (m) at (3.2,0) {\textbf{Mobile}};
\node (a) at (6.4,0) {\textbf{API}};
\node (d) at (9.6,0) {\textbf{DB}};
\draw[dashed] (u) -- +(0,-3.8);
\draw[dashed] (m) -- +(0,-3.8);
\draw[dashed] (a) -- +(0,-3.8);
\draw[dashed] (d) -- +(0,-3.8);
\draw[->] (0,-0.7) -- (3.2,-0.7) node[midway,above]{\tiny создаёт пост};
\draw[->] (3.2,-1.3) -- (6.4,-1.3) node[midway,above]{\tiny POST /feed};
\draw[->] (6.4,-1.9) -- (9.6,-1.9) node[midway,above]{\tiny INSERT feed\_posts};
\draw[<-] (6.4,-2.5) -- (9.6,-2.5) node[midway,above]{\tiny id,time};
\draw[<-] (3.2,-3.1) -- (6.4,-3.1) node[midway,above]{\tiny 201 + JSON};
\draw[<-] (0,-3.7) -- (3.2,-3.7) node[midway,above]{\tiny обновлённая лента};
\end{tikzpicture}
\caption{Sequence diagram: публикация поста в ленту}
\end{figure}

\subsection{ER-диаграмма (ядро данных)}

\begin{figure}[H]
\centering
\begin{tikzpicture}[node distance=9mm, >=Stealth]
\node[draw, rounded corners, fill=blue!8, minimum width=28mm] (users) {users};
\node[draw, rounded corners, fill=blue!8, below left=of users, minimum width=28mm] (pets) {pets};
\node[draw, rounded corners, fill=blue!8, below right=of users, minimum width=28mm] (feed) {feed\_posts};
\node[draw, rounded corners, fill=blue!8, right=of feed, minimum width=28mm] (comments) {feed\_comments};
\node[draw, rounded corners, fill=blue!8, right=of users, minimum width=28mm] (articles) {articles};
\draw[->] (pets) -- node[left]{\tiny owner\_id} (users);
\draw[->] (feed) -- node[left]{\tiny user\_id} (users);
\draw[->] (comments) -- node[above]{\tiny post\_id} (feed);
\draw[->] (comments) |- node[pos=0.75,right]{\tiny user\_id} (users);
\draw[->] (articles) -- node[above]{\tiny submitted\_by / approved\_by} (users);
\end{tikzpicture}
\caption{ER-диаграмма ключевых сущностей}
\end{figure}


\section{Технологический стек и обоснование выбора}

\subsection{Мобильная часть}

\textbf{React Native + Expo} выбраны как баланс между скоростью разработки и качеством пользовательского интерфейса:
\begin{itemize}
  \item единая кодовая база для iOS и Android;
  \item удобная интеграция с camera/location/image picker;
  \item быстрый цикл проверки изменений через Expo Go.
\end{itemize}

Альтернативы для мобильной части (\textit{Flutter}, нативная двойная разработка \textit{Swift+Kotlin}) рассматривались, но для данного проекта уступали по критерию скорости итераций и стоимости поддержки команды из одного разработчика. Приоритетом проекта был быстрый выпуск функционального вертикального среза (\textit{auth + map + social + moderation}) с минимальными overhead на инфраструктуру сборки.

\subsection{Серверная часть}

\textbf{Go + Gin} обеспечивают:
\begin{itemize}
  \item простой и быстрый HTTP-слой;
  \item понятную структуру handlers/services/router;
  \item устойчивую работу под параллельной нагрузкой.
\end{itemize}

В сравнении с \textit{Node.js/Express} и \textit{Django REST} выбор Go оправдан для проекта DogPaw следующими факторами: низкие накладные расходы на runtime, единый бинарный артефакт запуска, предсказуемая производительность для realtime-эндпоинтов карты и минимальная сложность деплоя учебного стенда.

\subsection{Хранилище}

\textbf{SQLite + GORM} подходят для учебного и пред-продакшен этапа:
\begin{itemize}
  \item быстрая инициализация без отдельного DB-сервера;
  \item автомиграции схемы;
  \item переносимость и низкий порог сопровождения.
\end{itemize}

Сравнение с альтернативами:
\begin{itemize}
  \item \textit{PostgreSQL} функционально богаче, но усложняет инфраструктуру и требует отдельного администрирования;
  \item \textit{MongoDB} упрощает документную модель, но усложняет связи и ролевые проверки в текущей предметной логике.
\end{itemize}
Для DogPaw выбран путь \textit{SQLite сейчас + возможность миграции на PostgreSQL}, что дает быстрый старт и сохраняет стратегию масштабирования.

\subsection{Карты и геоданные}

\textbf{react-native-maps} на клиенте и поиск мест на backend через \textbf{OpenStreetMap} (запросы к Overpass API) с кэшированием на сервере позволяют:
\begin{itemize}
  \item отображать маркеры по выбранным категориям и снижать лишние запросы при панорамировании (дебаунс, округление области);
  \item кэшировать повторные запросы к поиску мест на сервере в памяти на ограниченное время;
  \item периодически обновлять позиции пользователей на умной карте и при возврате на экран подтягивать актуальный список объявлений о пропаже.
\end{itemize}


\section{Реализация ключевых сценариев}

\subsection{Профиль питомца и QR}

При создании карточки формируется уникальный QR-идентификатор.  
Сценарии:
\begin{enumerate}
  \item пользователь может добавить фотографию питомца из галереи или сделать снимок с камеры прямо в приложении;
  \item при редактировании карточки доступно удаление ранее выбранного фото;
  \item владелец показывает QR на ошейнике;
  \item другой пользователь сканирует QR в приложении;
  \item система открывает карточку питомца и позволяет добавить владельца в друзья.
\end{enumerate}

\subsection{Прогулки}

Сценарий включает:
\begin{itemize}
  \item старт трекинга по разрешению геолокации;
  \item накопление дистанции, длительности (отображение с точностью до секунды) и оценки калорий; учёт расстояния с фильтрацией выбросов координат GPS;
  \item завершение и сохранение прогулки в историю;
  \item просмотр истории на отдельном экране и удаление записей.
\end{itemize}

\subsection{Лента}

Реализованы:
\begin{itemize}
  \item публикация текст + фото;
  \item лайк и его снятие;
  \item текстовые комментарии;
  \item сортировка по \textit{newest / likes / comments / interesting};
  \item удаление поста админом.
\end{itemize}

\subsection{Друзья}

Поддерживаются:
\begin{itemize}
  \item добавление по email;
  \item добавление через QR-поток;
  \item удаление из списка друзей.
\end{itemize}

\subsection{Умная карта и карта мест}

Реализовано:
\begin{itemize}
  \item фильтрация мест по четырем категориям (ветеринария, зоомагазин, груминг, парк/площадка);
  \item интерактивная легенда-контроллер;
  \item на карте мест отображается собственная геопозиция при выданном разрешении;
  \item видимость на умной карте: пользователь и друзья; скрытие друзей в статусе \code{do\_not\_disturb};
  \item цвет маркера соответствует статусу (в т.\,ч. ограничения платформы iOS для стандартных пинов);
  \item статус по умолчанию при обновлении позиции на сервере: \code{do\_not\_disturb}.
\end{itemize}

\subsection{База знаний и модерация}

Пользователь отправляет статью со статусом \code{pending}; просмотр опубликованной статьи выполняется на отдельном экране.  
Администратор видит отдельный блок ожидающих материалов, принимает решение:
\begin{itemize}
  \item \code{approve} $\rightarrow$ публикация;
  \item \code{reject} $\rightarrow$ отклонение;
  \item удаление записи (admin only).
\end{itemize}

\subsection{Потеря питомца}

Сценарий включает создание объявления из карточки питомца с произвольным текстом сообщения и координатами (по возможности --- текущее местоположение), отображение активных объявлений на карте в отдельном режиме и снятие объявления владельцем после нахождения питомца; список на клиенте обновляется при возврате на экран карты.


\section{REST API: спецификация и примеры}

\subsection{Карта endpoint'ов}

\begin{table}[H]
\centering
\small
\caption{Ключевые endpoint'ы backend API}
\begin{tabular}{|L{0.15\textwidth}|L{0.38\textwidth}|L{0.16\textwidth}|L{0.21\textwidth}|}
\hline
\textbf{Метод} & \textbf{Endpoint} & \textbf{Auth} & \textbf{Назначение} \\
\hline
POST & \code{/api/v1/auth/register} & нет & регистрация \\
\hline
POST & \code{/api/v1/auth/login} & нет & вход \\
\hline
GET & \code{/api/v1/auth/me} & JWT & профиль текущего пользователя \\
\hline
GET & \code{/api/v1/feed?sort=...} & нет & получить ленту \\
\hline
POST & \code{/api/v1/feed} & JWT & создать пост \\
\hline
POST & \code{/api/v1/feed/:id/comments} & JWT & добавить комментарий \\
\hline
POST & \code{/api/v1/friends} & JWT & добавить друга по id/email \\
\hline
DELETE & \code{/api/v1/friends/:id} & JWT & удалить друга \\
\hline
GET & \code{/api/v1/map/users} & JWT & умная карта: видимые пользователи \\
\hline
PUT & \code{/api/v1/map/me} & JWT & обновить свою позицию \\
\hline
POST & \code{/api/v1/articles/submit} & JWT & отправить статью \\
\hline
PUT & \code{/api/v1/articles/:id/moderate} & JWT(admin) & approve/reject \\
\hline
GET & \code{/api/v1/pets/:id/walks} & нет & история прогулок питомца \\
\hline
POST & \code{/api/v1/walks} & нет & сохранить прогулку \\
\hline
DELETE & \code{/api/v1/walks/:id} & JWT & удалить прогулку (владелец) \\
\hline
GET & \code{/api/v1/places/search?lat=...} & нет & поиск мест (внешние API + кэш) \\
\hline
GET/POST & \code{/api/v1/lost} & GET нет / POST JWT & объявления о пропаже \\
\hline
PUT & \code{/api/v1/lost/:id/found} & JWT & снять объявление (владелец) \\
\hline
\end{tabular}
\end{table}

\noindent\textit{Примечание:} в учебной сборке маршрут \code{POST /walks} не обёрнут в \code{authMiddleware} (в отличие от \code{DELETE /walks/:id}), что следует учитывать при переносе в production.

\subsection{Примеры запросов и JSON}

\textbf{1) Создание поста}
\begin{verbatim}
POST /api/v1/feed
Authorization: Bearer <token>
Content-Type: application/json

{
  "text": "Сегодня отличная прогулка!",
  "media_url": "data:image/jpeg;base64,..."
}
\end{verbatim}

\textbf{Ответ}
\begin{verbatim}
{
  "id": "post_uuid",
  "user_id": "user_uuid",
  "text": "Сегодня отличная прогулка!",
  "likes": 0,
  "created_at": "2026-04-01T09:15:00Z"
}
\end{verbatim}

\textbf{2) Отправка статьи на модерацию}
\begin{verbatim}
POST /api/v1/articles/submit
Authorization: Bearer <token>
Content-Type: application/json

{
  "title": "Как подготовить собаку к зиме",
  "content": "Пошаговый чеклист...",
  "category": "Уход"
}
\end{verbatim}

\textbf{Ответ}
\begin{verbatim}
{
  "id": "article_uuid",
  "status": "pending",
  "submitted_by": "user_uuid"
}
\end{verbatim}

\textbf{3) Обновление геопозиции}
\begin{verbatim}
PUT /api/v1/map/me
Authorization: Bearer <token>
Content-Type: application/json

{
  "latitude": 55.7558,
  "longitude": 37.6173,
  "status": "do_not_disturb"
}
\end{verbatim}

\noindent\textit{Примечание:} при отсутствии поля \code{status} сервер сохраняет значение по умолчанию \code{do\_not\_disturb}.

\subsection{Описание API (APU) и контрактов}

В проекте используется контрактный подход к API: клиентский модуль \code{src/api/client.ts} является единой точкой доступа и содержит типизированные методы для всех ключевых сценариев. На стороне backend маршруты объявлены централизованно в \code{router/router.go}, а проверки прав доступа вынесены в middleware (\code{authMiddleware}, \code{optionalAuthMiddleware}). Это снижает риск рассинхронизации между UI и серверной логикой.


\section{Тестирование и контроль качества}

\subsection{Стратегия тестирования}

Проверка выполнялась в несколько циклов:
\begin{enumerate}
  \item \textbf{Unit-level} проверки серверных пакетов (\code{go test ./...}) и типовой корректности фронтенда (\code{tsc});
  \item \textbf{Integration-level} smoke API сценариев: auth, feed, friends, map, articles, moderation;
  \item ручной E2E-прогон мобильных пользовательских потоков;
  \item ретест после исправлений;
  \item техническая валидация сборки и линтеров.
\end{enumerate}

Отдельный акцент сделан на интеграционные проверки ролевой логики: доступ к moderation-операциям, удаление контента администратором и owner-only выдача прав.

\subsection{Покрытые сценарии}

\begin{table}[H]
\centering
\small
\caption{Сценарии тестирования финального этапа}
\begin{tabular}{|L{0.5\textwidth}|L{0.14\textwidth}|L{0.26\textwidth}|}
\hline
\textbf{Сценарий} & \textbf{Статус} & \textbf{Комментарий} \\
\hline
Регистрация, вход, восстановление сессии & PASS & Проверены валидные и ошибочные ветки \\
\hline
Лента: пост, фото, лайк, комментарий, сортировки & PASS & Подтверждены все режимы сортировки \\
\hline
Друзья: add by email/QR, remove & PASS & Проверены симметричные проверки дружбы \\
\hline
Умная карта: видимость друзей и DND & PASS & DND корректно скрывает пользователя \\
\hline
Карта мест (поиск POI через OpenStreetMap/Overpass) & PASS & Проверена выдача и серверный кэш повторных запросов \\
\hline
База знаний: submit + moderation & PASS & Добавлен отдельный admin-блок pending \\
\hline
Owner-only выдача admin прав & PASS/CONDITIONAL & Проверен guard, полный e2e зависит от owner-логина \\
\hline
\end{tabular}
\end{table}

\subsection{Найденные проблемы и исправления}

Ключевые исправления по итогам тестирования:
\begin{itemize}
  \item устранены дубли backend/expo процессов, влияющие на непредсказуемость результата;
  \item переработан UX соц-экрана для уменьшения визуального шума;
  \item стабилизирован QR-сканер (защита от множественного срабатывания);
  \item исправлена обработка admin-контекста при чтении статей;
  \item настроена выдача мест через OpenStreetMap (Overpass) и кэширование повторных запросов на сервере;
  \item улучшена согласованность подписей статусов на умной карте с выбранным режимом.
\end{itemize}

\subsection{Документирование тестов}

Результаты по шагам зафиксированы в отдельном протоколе:
\begin{center}
\code{docs/ПРОТОКОЛ\_ТЕСТИРОВАНИЯ.md}
\end{center}
Это обеспечивает воспроизводимость QA-цикла и прозрачность ретестов.

\section{Визуальная концепция и UX-решения}

\subsection{Принципы UI}

\begin{itemize}
  \item \textbf{Иерархия:} заголовок $\rightarrow$ ключевые действия $\rightarrow$ контент.
  \item \textbf{Контраст:} текст и управляющие элементы читаемы при любых условиях.
  \item \textbf{Единообразие:} одинаковая логика кнопок и карточек на всех основных экранах.
  \item \textbf{Обратная связь:} явные статусы операций (успех, ошибка, режим кэша и т.д.).
\end{itemize}

\subsection{Пример UX-изменений}

\begin{enumerate}
  \item Карта мест: индикатор загрузки без блокировки всего экрана, снижение частоты запросов при движении карты.
  \item Энциклопедия: отдельный экран статьи; блок ожидающих модерации материалов для администратора.
  \item Социальный экран: снижен визуальный шум, улучшена читаемость полей ввода на тёмном фоне.
  \item Питомец и прогулки: единая тёмная тема; история прогулок и объявление о пропаже согласованы с картой.
\end{enumerate}

\subsection{Скриншоты приложения}

\begin{figure}[H]
\centering
\IfFileExists{screens/map.png}{
  \includegraphics[width=0.32\textwidth]{screens/map.png}
}{
  \fbox{\parbox[c][50mm][c]{0.30\textwidth}{\centering Скрин 1\\Карта мест}}
}
\hspace{4mm}
\IfFileExists{screens/social.png}{
  \includegraphics[width=0.32\textwidth]{screens/social.png}
}{
  \fbox{\parbox[c][50mm][c]{0.30\textwidth}{\centering Скрин 2\\Лента/друзья}}
}
\hspace{4mm}
\IfFileExists{screens/pet.png}{
  \includegraphics[width=0.32\textwidth]{screens/pet.png}
}{
  \fbox{\parbox[c][50mm][c]{0.30\textwidth}{\centering Скрин 3\\Карточка питомца}}
}
\caption{Ключевые экраны приложения DogPaw}
\end{figure}


\section{Риски и эксплуатация}

\subsection{Оценка трудозатрат}

Проект реализован итеративно в несколько этапов: архитектура, backend, frontend, интеграция, тестирование и документация. Основной вклад трудозатрат пришелся на:
\begin{itemize}
  \item проработку сквозных сценариев (auth $\leftrightarrow$ roles $\leftrightarrow$ content);
  \item отладку сетевого окружения мобильной разработки;
  \item синхронизацию контрактов API и экранной логики.
\end{itemize}

\subsection{Технические риски}

\begin{table}[H]
\centering
\small
\caption{Ключевые риски и меры снижения}
\begin{tabular}{|L{0.36\textwidth}|L{0.26\textwidth}|L{0.28\textwidth}|}
\hline
\textbf{Риск} & \textbf{Вероятность/влияние} & \textbf{Митигирующие меры} \\
\hline
Сетевые проблемы мобильного запуска в LAN & Средняя/Высокое & Переход между LAN/Tunnel, фиксированные порты, проверка firewall \\
\hline
Ошибки в ролевой модели доступа & Средняя/Высокое & Централизация middleware, отдельные e2e-ветки на admin/owner \\
\hline
Зависимость от внешнего API мест & Высокая/Среднее & OSM/Overpass; кэш на сервере; таймауты на клиенте; учёт лимитов публичных инстансов \\
\hline
Рост сложности UI при добавлении функций & Средняя/Среднее & UX-ревью, сокращение визуального шума, унификация компонентов \\
\hline
\end{tabular}
\end{table}

\subsection{План эксплуатации и развития}

При переходе к production-уровню целесообразны:
\begin{enumerate}
  \item перенос БД на PostgreSQL;
  \item выделенное файловое хранилище медиа;
  \item централизованный логинг и мониторинг;
  \item CI/CD с автотестами и проверками безопасности;
  \item публикация мобильных сборок через EAS и сторах;
  \item расширение медкарты питомца: история веса, план и журнал глистогонных и противоклещевых обработок;
  \item геймификация прогулок: достижения (самая длинная прогулка, серия дней) и дружеские челленджи/рейтинги.
\end{enumerate}


\section{Заключение}

В рамках финального этапа программного курсового проекта разработана и проверена интегрированная мобильная система DogPaw для владельцев собак. Проект прошел путь от прототипа до функционально связанного решения с авторизацией, социальной лентой, картографическим модулем, инструментами безопасности и модерацией пользовательского контента.

Главный результат состоит не только в наличии отдельных функций, но и в согласованности архитектуры: backend-контракты, правила ролей, клиентские сценарии и UX-слой работают как единая система. Отдельно подтверждена способность проекта к дальнейшему росту: текущая модульная структура позволяет масштабировать платформу, не нарушая базовую логику предметной области.


\begin{thebibliography}{99}

\bibitem{statista_pet}
Statista. Pet care market worldwide: statistics and trends, 2026. URL: \url{https://www.statista.com/topics/2849/pet-industry/}.

\bibitem{reactnative_docs}
React Native Documentation, 2026. URL: \url{https://reactnative.dev/docs/getting-started}.

\bibitem{expo_docs}
Expo Documentation, 2026. URL: \url{https://docs.expo.dev/}.

\bibitem{gin_docs}
Gin Web Framework Documentation, 2026. URL: \url{https://gin-gonic.com/docs/}.

\bibitem{gorm_docs}
GORM Documentation, 2026. URL: \url{https://gorm.io/docs/}.

\bibitem{owasp_api}
OWASP. API Security Top 10, 2023. URL: \url{https://owasp.org/API-Security/}.

\bibitem{sus}
Brooke J. SUS: A Quick and Dirty Usability Scale // Usability Evaluation in Industry. London: Taylor \& Francis, 1996.

\bibitem{fielding}
Fielding R. T. Architectural Styles and the Design of Network-based Software Architectures. Doctoral dissertation, University of California, Irvine, 2000.

\bibitem{expo_eas}
Expo. EAS Build and Submit Documentation, 2026. URL: \url{https://docs.expo.dev/build/introduction/}.

\bibitem{petdesk}
PetDesk. Official website, 2026. URL: \url{https://petdesk.com/}.

\bibitem{rover}
Rover. Official website, 2026. URL: \url{https://www.rover.com/}.

\bibitem{strava}
Strava. Official website, 2026. URL: \url{https://www.strava.com/}.

\bibitem{bringfido}
BringFido. Official website, 2026. URL: \url{https://www.bringfido.com/}.

\bibitem{pawboost}
PawBoost. Official website, 2026. URL: \url{https://www.pawboost.com/}.

\end{thebibliography}

\end{document}
